\scanpage{182}
en question induit sur \uline{Lie} est l'opération de puissance p.ème (est-ce dans le livre de Demazure-Gabriel ?). \margin{?} L'hom transposé est donc donné par l'opération de Cartier sur les formes différentielles.

\textcircled{7} \uline{Cas \inserted{lisse} de \deleted{\unsure{la}} dim relative un : un accouplement remarquable.}

\deleted{Soit} Plus généralement, au lieu de dim relative 1,
\[
 \deleted{Q=f_{X/S\,*}(\underline O_X)/\underline O_{X^{p/S}}}
\]
supposons que $X$ sur $X^{p/S}$ \deleted{\unsure{soit de loc une p-base,}} \inserted{soit de \unsure{t.f.|r.f.} et} $\Omega^1_{X/S}$ $(=\Omega^1_{X/X^{p/S}})$ soit \uline{libre de rang $1$}, ce qui signifie aussi que $X$ a loc. une $p$-base à 1 élément sur $S$. Posons
\[
\begin{aligned}
 Q&=f_{X/S\,*}(\underline O_X)/\operatorname{Im}\underline O_{X^{p/S}}
 =f_{X/S\,*}(\underline O_X)/\Lambda\\
 &\simeq\underline O_X/f^{-1}\underline O_S[\unsure{\underline o|\underline O}_X^p]
\end{aligned}
\]
qui est un module loc libre de rang $p-1$ sur $S$. On va définir un \uline{accouplement \inserted{$\Lambda$-bilinéaire} canonique}
\[
 \begin{aligned}
 Q\times Q&\longrightarrow\deleted{\unsure{\Lambda}}\,\Omega^1_{X^{p/S}/S}
 \otimes_{\underline O_{X^{p/S}}}\Lambda\\
 &\qquad(=\Omega^1_{X^{p/S}/S}\quad\text{si p.ex. $X/S$ lisse})
 \end{aligned}\tag*{($\ast$)}
\]
\margin{$-$} \deleted{ou} \deleted{\unsure{pour|plutôt}} ce qui revient au m\textsuperscript{e}, un hom $\Lambda$-bil.
\[
 \varphi:f_{X/S\,*}(\underline O_X)\times f_{X/S\,*}(\underline O_X)
 \longrightarrow\Omega^1_{X^{p/S}/S}\otimes_{\underline O_{X^{p/S}}}\Lambda
\]
qui est nul sur $\Lambda\times f_{X/S\,*}(\underline O_X)$ et sur $f_{X/S\,*}(\underline O_X)\times\Lambda$.
% archive sheet number: 181
